\begin{abstract}
Neural receivers (NRXs) replace channel estimation, equalization, and demapping with a trained neural network. A state-of-the-art NRX for the 5G NR uplink refines its estimates over iterations of convolutional and graph neural network (CGNN) layers, but within a 1\,ms latency budget, its real-time implementation admits only two of its eight iterations. Like other NRXs, it processes each slot independently and discards what earlier slots reveal about the channel. We propose a stateful NRX that keeps a pilot memory, i.e., the least-squares channel estimates at the pilot positions of earlier slots of each user, and combines it with the current pilots in a channel filter whose weights are not trained but fitted per user on the memory itself, without channel statistics or labels. In $2\times 4$ multi-user MIMO simulations with 16-QAM, the proposed two-iteration NRX requires more than 0.5\,dB less SNR than a stateless two-iteration NRX retrained on the same data, outperforms the three-iteration NRX, and stays within 0.2\,dB of the eight-iteration NRX, with and without power control. With QPSK, and on the tapped delay line and urban macrocell channels of two further datasets, it also outperforms the eight-iteration NRX. It meets the 1\,ms real-time budget, which the three-iteration NRX exceeds. With power control, combining the receiver's intermediate results with the cyclic redundancy check (CRC) history lets the base station request 97\% of the needed retransmissions before decoding finishes, against 79\% for the CRC history alone at the same false-alarm rate.
\end{abstract}
